\section*{Chapter \ref{ch:b3:hard-problems} --- Approximation and Exact Methods for Hard Problems}

\begin{solution}{exo:b3:hard-problems:1}
Edge $12$ joins the matching, $23$ is covered, $34$ joins, $45$ is covered, $56$ joins, $67$ is covered: the cover is $\{1, \ldots, 6\}$, six vertices. The smallest cover is $\{2, 4, 6\}$: three vertices, since the three disjoint edges $12$, $34$, $56$ need one each. Ratio $6/3 = 2$.
\end{solution}

\begin{solution}{exo:b3:hard-problems:2}
$S_1$ covers 4 new elements; then $S_4$ covers 5 and 6 (2 new, against 1 for $S_2$ and 1 for $S_3$). Two sets, $\{S_1, S_4\}$, and no single set covers $U$: optimal.
\end{solution}

\begin{solution}{exo:b3:hard-problems:3}
$x_1$ is a unit: true. Then $\neg x_1 \vee x_2$ is unit: $x_2$ true. Then $\neg x_2 \vee \neg x_3$: $x_3$ false. Then $x_3 \vee x_4 \vee \neg x_1$: $x_4$ true. Then $\neg x_4 \vee x_5$: $x_5$ true. Every variable is set, in the order $x_1, x_2, x_3, x_4, x_5$, and every clause holds.
\end{solution}

\begin{solution}{exo:b3:hard-problems:4}
Without the bound the count grows by about 15.5 for every four items, close to $2^4 = 16$: about 80 million nodes for 26 items. With the bound, the count depends on how close the fractional bound is to the true optimum on the instance, not only on its size; on the instance of 18 items it prunes more of the tree than on that of 14.
\end{solution}

\begin{solution}{exo:b3:hard-problems:5}
On the nine clauses of \code{test_exercise_5_buggy_is_unsound_and_fixed_is_right}, satisfiable, the buggy solver answers unsatisfiable. The first-UIP clause is $\neg u \vee \ell_1 \vee \cdots \vee \ell_k$, where the $\ell_i$ are literals of earlier levels; resolution derives it from the formula, so it is implied. Dropping the $\ell_i$ leaves the unit $\neg u$, which says $u$ is false whatever the earlier decisions; that is implied only if the conflict follows from $u$ alone, which is not the case when the earlier literals took part. The solver then asserts a consequence of its own decisions as a fact, and later refutes a satisfiable formula. The fix keeps the earlier literals and backjumps to the highest of their levels: \code{Fixed} uses the analysis of \cref{lst:b3:hard-problems:analyse} unchanged.
\end{solution}

\begin{solution}{exo:b3:hard-problems:6}
\code{exercises/ex6_one_tree.py} uses \code{minimum_spanning_tree} on the points but the first and adds the two shortest distances from the first. Any tour consists of the two tour edges at the first point, plus a path through all the others, which is a spanning tree of them: so the tour weighs at least the minimum spanning tree plus the two shortest edges at the first point (\code{test_tours}). For the van it gives 13,594~m.
\end{solution}

\begin{solution}{exo:b3:hard-problems:7}
Facts used: a matching's edges share no vertex; every cover contains an end of every edge.

Let $M$ be the matching built. An edge $uv$ not covered would have neither end in a matched edge, and the loop would have added it to $M$: so the result is a cover, of size $2|M|$. Any cover $C$ contains an end of each edge of $M$, and these ends are distinct, so $|C| \ge |M|$, and $2|M| \le 2|C|$. Tightness: on $k$ disjoint edges ($2k$ vertices) the algorithm takes all $2k$ vertices and $k$ suffice; on the complete bipartite graph of the chapter's figure, the same.
\end{solution}

\begin{solution}{exo:b3:hard-problems:8}
Facts used: a clause enters the queue only when its counter reaches 0, which happens once.

Building \code{waiting} and \code{uses} reads every literal once. Each clause is popped from the queue at most once, and its head is found by reading its literals: total, the size of the formula. A variable is made true at most once, and then each clause in which it occurs negatively has its counter decremented once: over all variables, at most the number of negative literals. Every step is constant time, so the whole is linear in the number of literals.
\end{solution}

\begin{solution}{exo:b3:hard-problems:9}
One entry per pair $(S, j)$ with $S \subseteq \{1, \ldots, n - 1\}$ and $j \in S$: $(n - 1)2^{n - 2}$ entries, each a minimum over at most $n - 2$ predecessors, so $O(n^2 2^n)$ steps. For $n = 12$: $11 \times 2^{10} = 11{,}264$ entries and at most about $11{,}264 \times 10 \approx 113{,}000$ steps (\code{test_van_bounds_and_table_size}), against $11!/2 = 19{,}958{,}400$ tours.
\end{solution}

\begin{solution}{exo:b3:hard-problems:10}
Facts used: $\sum_{j \le m} 1/j = H(m)$; greedy picks a set covering the most uncovered elements.

When greedy picks a set that covers $k$ new elements, charge each of them $1/k$; the total charge is the number of sets greedy takes. Let $S$ be a set of an optimal cover, and list its elements in the order greedy covers them, $e_1, \ldots, e_m$. When $e_i$ is covered, at least $m - i + 1$ elements of $S$ are still uncovered, so $S$ itself would cover that many; greedy's choice covers at least as many, and $e_i$ is charged at most $1/(m - i + 1)$. The elements of $S$ are charged at most $1/m + \cdots + 1/1 = H(m) \le H(n)$. Every element lies in some set of the optimal cover, so the total charge is at most $H(n)$ times the number of its sets. Finally $H(n) \le 1 + \int_1^n dx/x = 1 + \ln n$.
\end{solution}

\begin{solution}{exo:b3:hard-problems:11}
Facts used: \omterm{def:b3:np-completeness:problems}{Hamiltonian cycle} is \omterm{def:b3:np-completeness:reduction}{NP-complete} (\cref{ch:b3:np-completeness}).

Suppose an algorithm $A$ returns a tour within $\rho$ of the optimum in polynomial time. Given a graph $G$ on $n$ vertices, give each edge of $G$ length 1 and each other pair length $\rho n + 1$; this takes polynomial time. If $G$ has a \omterm{def:b3:np-completeness:problems}{Hamiltonian cycle}, the optimum is $n$ and $A$ returns a tour of length at most $\rho n$. If not, every tour uses a pair that is not an edge and has length at least $\rho n + 1 + (n - 1) > \rho n$. So comparing $A$'s tour with $\rho n$ decides \omterm{def:b3:np-completeness:problems}{Hamiltonian cycle}, and $\mathrm{P} = \mathrm{NP}$. The lengths break the triangle inequality, which is why the metric case escapes.
\end{solution}

\begin{solution}{exo:b3:hard-problems:12}
\code{exercises/ex12_restarts.py}: \code{Restarting} counts conflicts and, after $16 \cdot r_i$ of them, returns true from \code{restart_now}, so the main loop backjumps to level 0. The solver stays complete: the learned clauses are kept, each is implied by the formula so no solution is lost, and the run lengths grow without bound, so eventually a run is long enough to finish. Measured (\code{test_exercise_12_counts}): 119 conflicts without restarts against 90 with 4 restarts for 60 variables; 289 against 212 with 8 for 80; but 101 against 169 with 6 for 100. Restarts help on some formulas and hurt on others; solvers keep them because they cut the very long runs that an unlucky early decision causes.
\end{solution}

\begin{solution}{pb:b3:hard-problems:1}
\textbf{1.} $11!/2 = 19{,}958{,}400$.

\textbf{2.} Every tour passes through the Town Square; rotating it to start there gives the same cycle.

\textbf{3.} With $D(\{j\}, j) = d(0, j)$,
\[ D(S, j) = \min_{k \in S \setminus \{j\}} \bigl( D(S \setminus \{j\}, k) + d(k, j) \bigr), \]
and the tour is $\min_j D(\{1, \ldots, 11\}, j) + d(j, 0)$.

\textbf{4.} $11 \times 2^{10} = 11{,}264$.

\textbf{5.} About $10^5$ steps for 12 places (\cref{exo:b3:hard-problems:9}); for 24, $23 \times 2^{22} \approx 10^8$ entries of about 22 steps each, some 20,000 times more, minutes in Python, with a table too large for this program's dictionary.

\textbf{6.} No: Held--Karp is exact, so no tour is shorter than 16,635~m.

\textbf{7.} The straight line is the shortest path between two points, so going through a third point is never shorter.

\textbf{8.} Removing an edge from any tour leaves a spanning tree, which weighs at least the minimum one.

\textbf{9.} At most $2 \times 11{,}367 = 22{,}734$~m. Not met: 18,427~m, since the shortcuts save a great deal.

\textbf{10.} In replacing each skip of visited points by a straight segment.

\textbf{11.} $18{,}427 / 16{,}635 = 1.108$; the theorem promises at most 2.

\textbf{12.} The 1-tree, by $13{,}594 - 11{,}367 = 2{,}227$~m.

\textbf{13.} The optimum lies between 13,594 and 18,427~m, so \omterm{def:b3:hard-problems:tsp}{tree doubling} is within $18{,}427 / 13{,}594 = 1.36$ of it, a better guarantee for this instance than the factor 2.

\textbf{14.} Choose positions $i < j$ and reverse the places from $i$ to $j$; the edges into and out of the segment change.

\textbf{15.} Each move shortens the tour strictly, and there are finitely many tours.

\textbf{16.} It is a local optimum, since \code{two_opt} stops only when no move improves; it is not the optimum, 6\% longer.

\textbf{17.} Uncrossing two crossing edges by a \omterm{def:b3:hard-problems:local-search}{2-opt move} shortens the tour, as the two diagonals of a convex quadrilateral add up to more than either pair of opposite sides.

\textbf{18.} The pairs $i < j$: about $n^2/2$.

\textbf{19.} It starts within a factor 2, so fewer moves are needed, and the local optimum reached tends to be good; from a random tour 2-opt also works, but takes more moves.

\textbf{20.} Worst 1.456, mean 1.216.

\textbf{21.} Yes, $1.456 \le 2$. A ratio above 2 on a metric instance would contradict it, and so would show a bug.

\textbf{22.} After 2-opt the tours are on average 2.2\% above the optimum and at worst 15.6\%.

\textbf{23.} In 103 of 200: the other 97 stop at a local optimum that is not the best.

\textbf{24.} They cover 200 instances of one kind; the proof covers every metric instance, including those no sample finds.

\textbf{25.} Held--Karp would need about $60^2 \times 2^{60}$ steps: not feasible. \omterm{def:b3:hard-problems:tsp}{Tree doubling} and 2-opt run in a fraction of a second; the factor 2 remains, and the 1-tree bound says how close the tour is on that instance.
\end{solution}
